% Section 3 body fragment. Consent-procedure details should be completed from
% the ethics record before submission.
\section{Formative Study}
\label{sec:formative}

The formative study examined how participants reason about suspected coordinated
information operations (CIOs), what makes a coordination claim inspectable, and
how review support can preserve independent judgment. We used the study to
derive requirements for HyperTrace rather than to evaluate a finished
interface. We refer to this activity as a formative interview study, but the
archived analysis material is structured written response data rather than
audio-recorded conversational transcripts. The analysis package contains
anonymized responses from ten participants (P01--P10). The material analyzed for the
findings comprised the complete responses to questions Q3--Q6, yielding 40
response-level meaning units. Q2 was retained only as background information
and was not used to derive themes. Names, contact information, IP addresses,
and submission metadata were excluded before analysis.

\subsection{Participants and Materials}
\label{sec:formative-participants}

The ten participants had approximately two years of research experience on
average. The study was reviewed and approved by the institutional review board
(IRB) of the authors' college. We report these facts as study-context
information rather than inferring job titles or professional seniority from the
anonymous answers. This distinction is important because the themes describe
review practices and evidentiary needs, not the prevalence of those practices
in a particular occupational population.

The Q3--Q6 prompts elicited participants' accounts of (i) evidence used to
distinguish coordinated activity from coincidental or legitimate group
behavior, (ii) limitations and failure modes of model explanations, (iii) the
ability to trace an explanation to underlying records and make an independent
decision, and (iv) how explanations should change when new evidence arrives.
The prompts were framed as review scenarios rather than as a test of a
particular algorithm or interface. Quotations in this paper are translated
from the original Chinese responses; the source-language quote bank is retained
in the supplementary audit materials.

\begin{table}[t]
  \centering
  \caption{Participant and study-context information. The coding archive contains only the anonymous IDs and response material shown here.}
  \label{tab:formative-participants}
  \begin{tabular}{p{0.42\linewidth}p{0.45\linewidth}}
    \toprule
    Field & Status in current archive\\
    \midrule
    Participants & 10 (P01--P10)\\
    Analysis material & Structured written responses; Q3--Q6\\
    Meaning units & 40 complete response-level units\\
    Research experience & Approximately 2 years on average\\
    Recruitment channel & Not reported\\
    Ethics review & Approved by the authors' college IRB\\
    Consent procedure & To be completed from the ethics record\\
    \bottomrule
  \end{tabular}
\end{table}

\subsection{Coding and Analysis}
\label{sec:formative-analysis}

Two researchers independently performed AI-assisted open coding on the same
anonymized material and checked their respective coding outputs. After both
rounds, they compared the codes item by item. Synonymous labels were merged;
differences in coding granularity or theme boundaries were resolved through
consensus, with the decisions recorded in the analysis log. No contradictory
conclusion-level disagreement was identified. The analysis therefore combines
inductive attention to participants' wording with a deductive integration into
four high-level analytic targets (F1--F4). The second coding round retained ten
subthemes (T1--T10) to preserve within-theme distinctions and cross-cutting
issues.

We did not calculate support percentages or treat code frequency as evidence of
population prevalence. The unit of analysis was a complete response-level
meaning unit, and the purpose of coding was to identify recurring review
difficulties, evidence requirements, and risks that could inform system design.
The use of AI assistance is disclosed here and will be described in the
reproducibility materials; the researchers reviewed, compared, and documented
the resulting codes rather than presenting model-generated labels as
independent human coding.

\subsection{Findings}
\label{sec:formative-findings}

\subsubsection{F1: Coordination judgments require mechanism-level evidence, not risk aggregation}
\label{sec:f1}

All ten participants contributed a Q3 response coded to F1: P01, P02, P03,
P04, P05, P06, P07, P08, P09, and P10. T1 was represented by P01, P02, P06,
and P10; T2 by P03, P04, P05, and P07; P08 combined T1 and T2; and P09
combined T1 with the cross-cutting independent-judgment issue T8. These
identifiers refer to coded response units, not to a frequency estimate or a
claim that every participant endorsed every detail.

Participants distinguished a group of accounts that is jointly high-risk from
an operation that is coordinated through a shared mechanism. Content
similarity, similar risk scores, or a short burst of synchronized activity were
treated as clues rather than conclusions. Participants instead described
combining posting times, interaction targets, repeated behavioral patterns,
propagation paths, local graph structure, account life cycles, device or
account-asset links, and common generation templates. They also described
actively testing alternative explanations, including a news event, an interest
community, or a shared network service. As P01 explained:

\begin{quote}
\emph{``You cannot look only at whether these accounts' content is similar. You also need to consider posting times, retweet targets, interaction relationships, and behavioral patterns.''}
\end{quote}

The two subthemes separate behavioral regularity (T1) from shared relational,
account, or generation infrastructure (T2). This distinction prevents a
single salient signal from being treated as proof of coordination. It also
motivates evidence views that expose the relationships and event order behind
an alert, together with plausible benign alternatives.

\subsubsection{F2: Static prediction explanations do not cover adaptive coordination}
\label{sec:f2}

All ten participants contributed a Q4 response coded to F2: P01, P02, P03,
P04, P05, P06, P07, P08, P09, and P10. T3 was represented by P03, P04, and
P08; T4 by P01, P02, P06, and P10; P07 and P09 contributed T5; and P05's
response was mapped to the cross-cutting T7 boundary concerning evidence-chain
accountability. The T7 mapping is retained in the codebook because it connects
explanation failure to downstream evidentiary responsibility; it does not
change F2's main claim.

Participants viewed a single post explanation, local feature attribution, or
static salient subgraph as insufficient for explaining how a multi-account
operation formed or whether the explanation remained valid at the relevant
time. They anticipated explanation failures caused by adversarial noise,
structural camouflage, distribution shift, missing social context, implicit
meaning, slow generation, and excessive information density. P08 summarized the
temporal problem as follows:

\begin{quote}
\emph{``Even if XAI provides an explanation, it may be using old data patterns to explain a new attack method.''}
\end{quote}

The subthemes capture both adversarial or adaptive change (T3) and the limits
of local attribution for compositional and socially situated behavior (T4).
Engineering and cognitive usability cut across these concerns (T5): an
explanation that is slow, overly technical, or simultaneously presents too
many graphs and scores may be less usable even when its underlying attribution
is technically correct. The boundary of this finding is equally important:
participants did not claim that every static explanation is useless; they
identified conditions under which it can become stale, incomplete, or
overwhelming.

\subsubsection{F3: Explanation claims must be traceable item by item to source records}
\label{sec:f3}

All ten participants contributed a Q5 response coded to F3: P01, P02, P03,
P04, P05, P06, P07, P08, P09, and P10. T6 was represented by P01, P02, P03,
P04, P05, P06, P07, P08, and P09; T7 by P05 and P10; and T8 by P01, P02,
and P10. P07 and P09 were also mapped to the cross-cutting usability theme
T5. Participants may therefore contribute to more than one subtheme within
the same response.

Participants required each important explanatory claim to resolve to the
corresponding account, content item, relation, timestamp, or behavioral log.
Source records were needed not only to establish that an item existed, but also
to check relation direction, surrounding context, and whether the item was
available at the decision time. P03 stated:

\begin{quote}
\emph{``Every node and edge in the explanation should be traceable back to the original data.''}
\end{quote}

T6 concerns technical reachability from an explanation element to a source
record. T7 extends this requirement to chain-of-custody, appeal, regulatory,
or other externally accountable settings. T8 emphasizes that source access is
also a human-decision requirement: participants rejected treating an AI label
as evidence in itself. A reviewer should be able to inspect the underlying
record, qualify the claim, and retain the option to disagree. Participants also
mentioned counterfactual checks, but the analysis distinguishes a model's
counterfactual dependence from proof of real-world causality or legal evidence.

\subsubsection{F4: Dynamic review requires explicit and reversible explanation updates}
\label{sec:f4}

All ten participants contributed a Q6 response coded to F4: P01, P02, P03,
P04, P05, P06, P07, P08, P09, and P10. T9 was represented by P01, P02, P03,
P05, P06, and P09; T10 by P02, P04, P07, P08, P09, and P10. P02 and P09
also contributed the cross-cutting T5 concern about reminder load and
interruption control. The overlap reflects the codebook's mapping of one
response to multiple related subthemes, not additional participants.

Participants wanted new evidence to be presented as an explicit increment to an
existing assessment rather than silently replacing the previous explanation.
They asked to see which events, accounts, relations, risk states, and
uncertainties changed, why the change occurred, and what evidence supported the
new state. P05 described the need for a persistent history:

\begin{quote}
\emph{``In addition to showing the current conclusion, the system should retain the previous judgment, the rules and evidence used at that time, and why it later changed.''}
\end{quote}

T9 therefore concerns version history, graph differences, historical
snapshots, and rollback. T10 concerns changes in risk level, confidence or
uncertainty, threshold escalation, and task routing. Participants did not ask
for an alert on every new event. Instead, reminders should be tied to a
substantive change in risk, evidence, or uncertainty and should remain open to
human review; otherwise, item-level notifications may create alert fatigue.

\subsection{From Findings to Design Requirements}
\label{sec:formative-requirements}

We translated the four findings into four requirements for HyperTrace. The
requirements describe what the review workflow must make possible; they do not
claim that the interview participants specified the implementation algorithm.
The mapping is summarized in Table~\ref{tab:formative-requirements}.

\begin{table*}[t]
  \centering
  \caption{Mapping formative findings to HyperTrace design requirements and later evaluation.}
  \label{tab:formative-requirements}
  \begin{tabular}{p{0.18\linewidth}p{0.25\linewidth}p{0.30\linewidth}p{0.18\linewidth}}
    \toprule
    Finding & Design requirement & HyperTrace response & Evaluation link\\
    \midrule
    F1: mechanism-level coordination evidence & DR1: expose converging behavioral, relational, and temporal signals and preserve alternative interpretations & Typed evidence subgraph, event sequence, shared targets, and mechanism-level evidence labels & Evidence verification and final decision accuracy\\
    F2: static explanations can become stale or overload reviewers & DR2: make scope, uncertainty, and explanation limits visible without overwhelming the reviewer & Cutoff and scope state, compact evidence budget, warnings, and progressive disclosure & Temporal validity, sparsity, stability, and workload\\
    F3: claims must reach original records & DR3: provide item-level evidence provenance and support independent challenge & Evidence IDs, source records, timestamps, relation direction, and source-verification actions & Provenance coverage and source-verification accuracy\\
    F4: updates must be explicit and reversible & DR4: preserve changes, reasons, and uncertainty across review states & Version history, graph diff, changed margins, reminders, and rollback/review actions & Update traceability and erroneous-recommendation rejection\\
    \bottomrule
  \end{tabular}
\end{table*}

\begin{figure*}[t]
  \centering
  \fbox{\parbox{0.92\textwidth}{\centering
    \vspace{0.25in}
    \textbf{Figure placeholder: formative findings to design requirements.}\\
    F1--F4 and T1--T10 should map to DR1--DR4, HyperTrace interface elements,
    and the corresponding technical or human evaluation measures.\\
    \vspace{0.25in}
  }}
  \caption{Formative findings translated into design requirements, HyperTrace responses, and evaluation measures. The final figure should use only the themes and mappings reported in this section.}
  \label{fig:formative-to-requirements}
\end{figure*}

\subsection{Methodological Transparency and Scope}
\label{sec:formative-transparency}

The study provides design evidence, not a prevalence estimate or a claim that
all platform reviewers reason in the same way. Its sample, response format,
language, and the researchers' relationship to the project should be reported
in the final submission. We will also report the
translation procedure for the English quotations and retain a de-identified
question guide, codebook, theme map, quote bank, comparison notes, and audit
log in the supplementary materials where permitted. Because the coding used
AI assistance, the paper will state what material was provided to the coding
tool, how identifying information was removed, what researchers checked, and
which decisions were made by consensus. No reliability statistic is reported:
the current archive documents independent AI-assisted coding, comparison, and
consensus merging, but not a completed inter-rater reliability calculation.

These boundaries shape the rest of the paper. HyperTrace is evaluated against
the needs identified here, but the formative study itself does not establish
that the system improves reviewer accuracy. That question is addressed
separately through the technical explanation evaluation and the controlled
human review study.
